\documentclass[aspectratio=169,11pt]{beamer}

\usetheme{Madrid}
\usecolortheme{default}

\usepackage[T1]{fontenc}
\usepackage[utf8]{inputenc}
\usepackage{booktabs}
\usepackage{tabularx}
\usepackage{array}
\usepackage{xurl}

\title[Thesis Progress]{Evidence-Based Developer Reputation for Open-Source Software}
\subtitle{Progress presentation: from dataset construction to metric selection evidence}
\author{Ali Bakhshayesh}
\institute{University of Padova}
\date{\today}

\newcommand{\metric}[1]{\textbf{#1}}
\newcommand{\smalltext}[1]{{\footnotesize #1}}

\begin{document}

\begin{frame}
  \titlepage
\end{frame}

\begin{frame}{Purpose of This Presentation}
  \begin{itemize}
    \item Explain the thesis idea in simple language.
    \item Show what has been built so far, starting from zero.
    \item Explain how the dataset was created and cleaned.
    \item Explain the candidate metrics and comparison criteria.
    \item Show the current result: which metric is strongest for the PoC.
    \item Clarify what still needs to be done next.
  \end{itemize}
\end{frame}

\begin{frame}{The Big Picture}
  The thesis is about a simple question:

  \vspace{0.4cm}
  \begin{block}{Main question}
    Can we create a trustworthy developer-reputation signal from open-source software evidence?
  \end{block}

  \vspace{0.3cm}
  In other words, we want to measure developer activity and responsibility using evidence that can be checked again later.
\end{frame}

\begin{frame}{Why Developer Reputation Matters}
  Open-source projects often depend on many people:
  \begin{itemize}
    \item some people make small occasional changes;
    \item some people maintain the project for a long time;
    \item some people do large technical work;
    \item some activity may come from bots or automation.
  \end{itemize}

  \vspace{0.3cm}
  The challenge is to identify useful signals without making unfair or uncheckable claims.
\end{frame}

\begin{frame}{The DESTINO Context}
  The thesis is connected to a DESTINO-style evidence workflow.

  \vspace{0.2cm}
  The idea is:
  \begin{enumerate}
    \item a developer submits software evidence;
    \item the evidence is linked to a hash, so it cannot silently change;
    \item metrics are computed from that evidence;
    \item compact results are stored in a smart-contract style record;
    \item another person can later replay the measurement and check the result.
  \end{enumerate}
\end{frame}

\begin{frame}{Important Design Choice}
  We do \textbf{not} try to compute repository metrics directly inside a smart contract.

  \vspace{0.3cm}
  That would be too expensive and too complex because Git history analysis needs:
  \begin{itemize}
    \item repository traversal;
    \item file and line counting;
    \item commit and author analysis;
    \item bot filtering;
    \item normalization and validation.
  \end{itemize}

  \vspace{0.3cm}
  Instead, we use a \textbf{hybrid design}: compute off-chain, store compact auditable results on-chain.
\end{frame}

\begin{frame}{The Thesis Workflow}
  The work is organized in two main parts.

  \vspace{0.3cm}
  \begin{block}{Phase I: Build the evidence base}
    Select repositories, freeze versions, build release windows, extract metrics, clean data, and validate outputs.
  \end{block}

  \begin{block}{Phase II: Compare candidate metrics}
    Evaluate which metric is stable, understandable, resistant to gaming, useful long-term, and feasible for the PoC.
  \end{block}
\end{frame}

\begin{frame}{Where We Started}
  We first tested the workflow on a small pilot dataset.

  \begin{itemize}
    \item The pilot used 4 repositories.
    \item The goal was to understand the extraction flow before scaling.
    \item After the flow worked, we expanded the dataset.
  \end{itemize}

  \vspace{0.3cm}
  This reduced risk: we tested the method first, then applied it to more projects.
\end{frame}

\begin{frame}{Dataset Expansion}
  The dataset was expanded from the initial 4 repositories to 10 repositories.

  \vspace{0.2cm}
  Current repositories:
  \begin{itemize}
    \item \texttt{brew}, \texttt{click}, \texttt{curl}, \texttt{flask}, \texttt{gin}
    \item \texttt{junit5}, \texttt{prettier}, \texttt{requests}, \texttt{tokio}, \texttt{urllib3}
  \end{itemize}

  \vspace{0.3cm}
  All evidence is now centralized. We no longer depend on separate pilot-only or extended-only CSV files.
\end{frame}

\begin{frame}{What ``Dataset Freeze'' Means}
  A dataset freeze means that we fixed the input projects before analysis.

  \vspace{0.2cm}
  For each repository we recorded:
  \begin{itemize}
    \item project identifier;
    \item repository URL;
    \item persistence status;
    \item software-type tags;
    \item inclusion metadata.
  \end{itemize}

  \vspace{0.3cm}
  This is important because later results must be tied to a clear and fixed dataset.
\end{frame}

\begin{frame}{What ``Version Freeze'' Means}
  For each repository, we selected the latest 8 stable semantic-version tags.

  \vspace{0.2cm}
  We excluded unstable release labels such as:
  \begin{itemize}
    \item alpha;
    \item beta;
    \item release candidate versions.
  \end{itemize}

  \vspace{0.3cm}
  Each selected version is tied to an immutable commit hash. This means the exact code state can be checked again later.
\end{frame}

\begin{frame}{Release Windows}
  After freezing versions, we built release windows.

  \vspace{0.2cm}
  A release window is the interval between two consecutive tags.

  \vspace{0.3cm}
  Example:
  \begin{block}{Simple example}
    Version 1.0.0 to Version 1.1.0 is one window.
  \end{block}

  \vspace{0.3cm}
  We have:
  \begin{itemize}
    \item 10 repositories;
    \item 8 frozen tags per repository;
    \item 7 windows per repository;
    \item 70 total release windows.
  \end{itemize}
\end{frame}

\begin{frame}{Why Release Windows Matter}
  We do not only want to know what a project looks like at one moment.

  \vspace{0.2cm}
  We want to see what changed between two releases:
  \begin{itemize}
    \item how many commits happened;
    \item how much code changed;
    \item who contributed;
    \item whether ownership was concentrated in one or a few developers.
  \end{itemize}

  \vspace{0.3cm}
  This makes the analysis more useful for developer reputation.
\end{frame}

\begin{frame}{Candidate Metrics}
  We studied five candidate metric families.

  \begin{table}
    \footnotesize
    \begin{tabularx}{\textwidth}{p{1.0cm}p{3.7cm}X}
      \toprule
      ID & Metric & Simple meaning \\
      \midrule
      M1 & Snapshot size & How large the project is at a release tag. \\
      M2 & Churn & How many lines were added or deleted in a release window. \\
      M3 & Activity cadence & How often developers are active through commits and active days. \\
      M4 & Ownership/concentration & How much of the work is done by top contributors. \\
      M5 & Responsiveness & How quickly maintainers respond to issues or pull requests. \\
      \bottomrule
    \end{tabularx}
  \end{table}
\end{frame}

\begin{frame}{Important Note About M5}
  M5, maintenance responsiveness, is useful conceptually.

  \vspace{0.3cm}
  But it needs issue and pull-request event data.

  \vspace{0.3cm}
  Our current Phase I dataset focuses on:
  \begin{itemize}
    \item Git tags;
    \item commits;
    \item code changes;
    \item developer identities in Git history.
  \end{itemize}

  \vspace{0.3cm}
  Therefore, M5 is kept as a future extension, not as a current PoC candidate.
\end{frame}

\begin{frame}{Metric Extraction}
  We extracted four implemented metric families from the frozen repositories:

  \begin{itemize}
    \item \metric{M1:} lines of code and source lines of code per frozen tag.
    \item \metric{M2:} added lines, deleted lines, and total churn per window.
    \item \metric{M3:} commits, active days, and activity rate.
    \item \metric{M4:} ownership shares, top contributor shares, and Gini concentration.
  \end{itemize}

  \vspace{0.3cm}
  These outputs form the empirical base for the later comparison.
\end{frame}

\begin{frame}{Cleaning Developer Identity}
  Developer identity can be messy.

  \vspace{0.2cm}
  The same person may appear with different names, but the same email.

  \vspace{0.3cm}
  We canonicalized developer identity mainly by email. This makes developer-level metrics more consistent.

  \vspace{0.3cm}
  Cleaning result:
  \begin{itemize}
    \item clean rows remained stable;
    \item duplicate groups merged: 0;
    \item zero-commit windows: 0;
    \item zero-churn windows: 0;
    \item error log: empty.
  \end{itemize}
\end{frame}

\begin{frame}{Bot Handling}
  Some repository activity comes from bots or automated accounts.

  \vspace{0.2cm}
  We did not silently mix bot activity with human developer behavior.

  \vspace{0.3cm}
  Instead:
  \begin{itemize}
    \item bot rows were flagged;
    \item bot summaries were kept for audit;
    \item main M3 and M4 evidence uses human-only outputs.
  \end{itemize}

  \vspace{0.3cm}
  Bots were 15.85\% of commits but only 1.61\% of churn lines.
\end{frame}

\begin{frame}{Duration Normalization}
  Release windows are not all the same length.

  \vspace{0.2cm}
  One project may release after 3 days, another after 6 months.

  \vspace{0.3cm}
  To compare fairly, we computed normalized values:
  \begin{itemize}
    \item commits per day;
    \item churn per day;
    \item churn per KLOC per day.
  \end{itemize}

  \vspace{0.3cm}
  Very short windows are kept but flagged for sensitivity analysis.
\end{frame}

\begin{frame}{Centralized Evidence}
  All important Phase I evidence is now centralized in \texttt{\_all} files.

  \vspace{0.2cm}
  Examples:
  \begin{itemize}
    \item \texttt{dataset\_freeze\_all.csv}
    \item \texttt{versions\_all.csv}
    \item \texttt{windows\_all.csv}
    \item \texttt{metrics\_window\_dev\_human\_all.csv}
    \item \texttt{ownership\_summary\_human\_all.csv}
  \end{itemize}

  \vspace{0.3cm}
  This makes the project easier to reproduce and avoids duplicate parallel files.
\end{frame}

\begin{frame}{Phase I Result}
  Phase I produced a clean, reproducible evidence base.

  \vspace{0.3cm}
  In simple terms:
  \begin{itemize}
    \item we know exactly which repositories are included;
    \item we know exactly which versions were analyzed;
    \item we know exactly which release windows were measured;
    \item we extracted and cleaned the metric evidence;
    \item we separated human activity from bot activity;
    \item we normalized values so windows can be compared more fairly.
  \end{itemize}
\end{frame}

\begin{frame}{Why Phase II Was Needed}
  After Phase I, we had metrics.

  \vspace{0.3cm}
  But we still needed to answer:
  \begin{block}{Key question}
    Which metric is best for a developer-reputation PoC?
  \end{block}

  \vspace{0.3cm}
  A metric is not good just because it is easy to compute. It also needs to be stable, understandable, hard to manipulate, useful over time, and feasible for the target system.
\end{frame}

\begin{frame}{Phase II Criteria}
  We compared the candidate metrics using four criteria.

  \begin{table}
    \footnotesize
    \begin{tabularx}{\textwidth}{p{1.0cm}p{3.5cm}X}
      \toprule
      ID & Criterion & Simple question \\
      \midrule
      C1 & Stability & Does the metric behave consistently over time? \\
      C2 & Interpretability & Can we explain the metric clearly? \\
      C3 & Gaming resistance & Is the metric hard to manipulate? \\
      C4 & Long-term suitability & Does it stay useful as history grows? \\
      \bottomrule
    \end{tabularx}
  \end{table}
\end{frame}

\begin{frame}{Clarification: C1-C4 vs Feasibility}
  There are two different evaluations.

  \vspace{0.2cm}
  \begin{block}{C1-C4}
    These ask whether a metric is good as a reputation signal.
  \end{block}

  \begin{block}{Smart-contract feasibility}
    This asks whether the metric can be implemented and audited safely in the PoC system.
  \end{block}

  \vspace{0.2cm}
  C2, C3, and C4 each use five local dimensions. The feasibility rubric is separate and uses ten dimensions.
\end{frame}

\begin{frame}{C1: Stability}
  C1 checked how stable the metrics are across release windows.

  \vspace{0.2cm}
  Main findings:
  \begin{itemize}
    \item M1 is very stable, but it is mostly project context.
    \item M2 churn is informative but volatile.
    \item M3 activity has useful rank continuity among recurring developers.
    \item M4 concentration signals are more stable than raw activity.
  \end{itemize}

  \vspace{0.3cm}
  C1 did not select the final metric. It produced the first evidence layer.
\end{frame}

\begin{frame}{C2: Interpretability}
  C2 checked whether each metric is easy to explain.

  \begin{table}
    \footnotesize
    \begin{tabular}{lcc}
      \toprule
      Metric & Score & Band \\
      \midrule
      M3 Activity cadence & 9/10 & High \\
      M1 Snapshot size & 8/10 & High \\
      M2 Churn & 8/10 & High \\
      M4 Ownership/concentration & 7/10 & Medium \\
      M5 Responsiveness & 6/10 & Medium \\
      \bottomrule
    \end{tabular}
  \end{table}

  \vspace{0.2cm}
  Main message: M3 is easiest to explain, while M4 is meaningful but needs more explanation.
\end{frame}

\begin{frame}{C3: Gaming Resistance}
  C3 checked how easy each metric is to manipulate.

  \begin{table}
    \footnotesize
    \begin{tabular}{lcc}
      \toprule
      Metric & Score & Band \\
      \midrule
      M4 Ownership/concentration & 9/10 & High \\
      M2 Churn & 7/10 & Medium \\
      M3 Activity cadence & 6/10 & Medium \\
      M1 Snapshot size & 5/10 & Medium \\
      M5 Responsiveness & 3/10 & Low \\
      \bottomrule
    \end{tabular}
  \end{table}

  \vspace{0.2cm}
  Main message: M4 is strongest because sustained ownership is harder to fake than simple commits or churn.
\end{frame}

\begin{frame}{C4: Long-Term Suitability}
  C4 checked whether the metric remains useful as project history grows.

  \begin{table}
    \footnotesize
    \begin{tabular}{lcc}
      \toprule
      Metric & Score & Band \\
      \midrule
      M3 Activity cadence & 8/10 & High \\
      M4 Ownership/concentration & 8/10 & High \\
      M1 Snapshot size & 7/10 & Medium \\
      M2 Churn & 7/10 & Medium \\
      M5 Responsiveness & 4/10 & Low \\
      \bottomrule
    \end{tabular}
  \end{table}

  \vspace{0.2cm}
  Main message: M3 and M4 both work well over longer histories, but they mean different things.
\end{frame}

\begin{frame}{Combined C1-C4 Result}
  After C1-C4, the strongest implemented reputation-oriented metric is M4.

  \begin{table}
    \footnotesize
    \begin{tabular}{lcc}
      \toprule
      Metric & C1-C4 average & Role \\
      \midrule
      M4 Ownership/concentration & 8.13 & Strongest candidate \\
      M1 Snapshot size & 7.50 & Context only \\
      M3 Activity cadence & 7.25 & Eligible alternative \\
      M2 Churn & 6.50 & Supporting signal \\
      M5 Responsiveness & 3.25 & Future work \\
      \bottomrule
    \end{tabular}
  \end{table}
\end{frame}

\begin{frame}{Smart-Contract Feasibility}
  After comparing the metrics scientifically, we checked implementation feasibility.

  \vspace{0.2cm}
  The feasibility rubric asks whether a metric can be:
  \begin{itemize}
    \item computed deterministically;
    \item stored compactly;
    \item updated at a realistic frequency;
    \item audited from hash-pinned evidence;
    \item represented safely in a smart-contract record;
    \item protected against simple manipulation.
  \end{itemize}

  \vspace{0.2cm}
  This rubric has 10 dimensions, D1-D10.
\end{frame}

\begin{frame}{Feasibility Results}
  \begin{table}
    \footnotesize
    \begin{tabular}{lccc}
      \toprule
      Metric & Score & Band & Status \\
      \midrule
      M1 Snapshot size & 18/20 & High & Context only \\
      M2 Churn & 16/20 & High & Eligible \\
      M3 Activity cadence & 17/20 & High & Eligible \\
      M4 Ownership/concentration & 16/20 & High & Eligible \\
      M5 Responsiveness & 10/20 & Medium & Excluded \\
      \bottomrule
    \end{tabular}
  \end{table}

  \vspace{0.2cm}
  M1 is technically feasible, but it is not selected because size is not a developer-reputation metric by itself.
\end{frame}

\begin{frame}{Final Decision Inputs}
  The final decision input combines empirical quality and feasibility.

  \begin{table}
    \footnotesize
    \begin{tabular}{lccc}
      \toprule
      Metric & C1-C4 avg. & Feasibility & Combined \\
      \midrule
      M4 Ownership/concentration & 8.13 & 8.0 & 8.08 \\
      M3 Activity cadence & 7.25 & 8.5 & 7.75 \\
      M2 Churn & 6.50 & 8.0 & 7.10 \\
      M1 Snapshot size & 7.50 & 9.0 & Context only \\
      M5 Responsiveness & 3.25 & 5.0 & Excluded \\
      \bottomrule
    \end{tabular}
  \end{table}

  \vspace{0.2cm}
  Current result: M4 is preferred for the final selection gate.
\end{frame}

\begin{frame}{Why M4 Is Currently the Best Candidate}
  M4 measures ownership and concentration.

  \vspace{0.2cm}
  It is strong because:
  \begin{itemize}
    \item it is more reputation-oriented than project size;
    \item it is less volatile than churn alone;
    \item it is harder to game than simple commit counts;
    \item it remains useful over longer histories;
    \item it can be implemented in the hybrid PoC model.
  \end{itemize}

  \vspace{0.2cm}
  Main caveat: identity handling must be explicit and auditable.
\end{frame}

\begin{frame}{How M2 and M3 Still Matter}
  M4 should not be interpreted alone.

  \vspace{0.2cm}
  Supporting signals:
  \begin{itemize}
    \item \metric{M2 churn} checks whether ownership is backed by real code change.
    \item \metric{M3 activity cadence} checks whether the developer is active over time.
    \item \metric{M1 size} gives project context and normalization support.
  \end{itemize}

  \vspace{0.3cm}
  The PoC may select one primary metric, but the analysis keeps the supporting evidence visible.
\end{frame}

\begin{frame}{What Has Been Committed}
  The project is now versioned through Git.

  \vspace{0.2cm}
  Recent committed milestones:
  \begin{itemize}
    \item centralized Phase I ten-repository evidence;
    \item C1 stability analysis;
    \item C2 interpretability analysis;
    \item C3 gaming-resistance analysis;
    \item C4 long-term suitability analysis;
    \item smart-contract feasibility scoring.
  \end{itemize}

  \vspace{0.2cm}
  The branch is ahead of the remote repository because these local commits have not been pushed yet.
\end{frame}

\begin{frame}{What Is New After the Last Clarification}
  The report now explicitly explains that:

  \begin{itemize}
    \item C2, C3, and C4 each use five criterion-specific dimensions.
    \item The feasibility rubric is separate and uses ten dimensions.
    \item These are two different computations with different purposes.
  \end{itemize}

  \vspace{0.3cm}
  This was added to prevent confusion between \texttt{C2\_D1--C2\_D5} and feasibility \texttt{D1--D10}.
\end{frame}

\begin{frame}{Current State}
  At this point:
  \begin{itemize}
    \item Phase I dataset construction is complete.
    \item Metrics M1-M4 were extracted and cleaned.
    \item Bot handling and duration normalization are documented.
    \item Phase II C1-C4 comparison is complete.
    \item Smart-contract feasibility scoring is complete.
    \item M4 is the preferred candidate for the final PoC selection gate.
  \end{itemize}
\end{frame}

\begin{frame}{Next Step}
  The next step is not another metric extraction step.

  \vspace{0.3cm}
  The next step is:
  \begin{block}{Final metric selection rationale}
    Officially record why M4 is selected as the baseline PoC metric, and define exactly how it will be represented in the PoC.
  \end{block}

  \vspace{0.3cm}
  This should freeze:
  \begin{itemize}
    \item selected metric definition;
    \item analysis policy;
    \item value encoding;
    \item evidence hash meaning;
    \item smart-contract record fields.
  \end{itemize}
\end{frame}

\begin{frame}{After Metric Selection}
  After the final metric is selected, we can start the PoC implementation.

  \vspace{0.2cm}
  The PoC should include:
  \begin{itemize}
    \item a minimal smart-contract record structure;
    \item a compact metric value;
    \item an evidence hash;
    \item a submission or registration flow;
    \item an off-chain verifier that computes the selected metric;
    \item one or more test examples from the frozen dataset.
  \end{itemize}
\end{frame}

\begin{frame}{Summary}
  \begin{itemize}
    \item We built a clean ten-repository evidence base.
    \item We froze versions, windows, and commit hashes for reproducibility.
    \item We extracted and cleaned M1-M4 metric evidence.
    \item We evaluated the metrics using C1-C4.
    \item We applied the smart-contract feasibility rubric.
    \item The current strongest PoC candidate is M4 ownership/concentration.
    \item The next step is to write the final selection rationale and then implement the PoC.
  \end{itemize}
\end{frame}

\begin{frame}{One-Sentence Conclusion}
  \begin{block}{Conclusion}
    The project has moved from raw repository selection to a reproducible evidence base and a justified PoC candidate, with M4 ownership/concentration currently preferred because it gives the best tradeoff between reputation meaning, robustness, long-term usefulness, and implementation feasibility.
  \end{block}
\end{frame}

\end{document}
